
\documentclass{article}
\usepackage{colortbl}
\usepackage{makecell}
\usepackage{multirow}
\usepackage{supertabular}

\begin{document}

\newcounter{utterance}

\centering \large Interaction Transcript for game `function\_detective', experiment `string\_passive\_examples', episode 163 with mistral{-}large{-}3{-}azure.
\vspace{24pt}

{ \footnotesize  \setcounter{utterance}{1}
\setlength{\tabcolsep}{0pt}
\begin{supertabular}{c@{$\;$}|p{.15\linewidth}@{}p{.15\linewidth}p{.15\linewidth}p{.15\linewidth}p{.15\linewidth}p{.15\linewidth}}
    \# & \multicolumn{2}{c}{Player} && \multicolumn{2}{c}{Game Master} \\
    \hline

    \theutterance \stepcounter{utterance}  
    & & & \multicolumn{4}{p{0.6\linewidth}}{
        \cellcolor[rgb]{0.9,0.9,0.9}{
            \makecell[{{p{\linewidth}}}]{
                \texttt{\tiny{[P1$\langle$GM]}}
                \texttt{I have a hidden Python function for you to figure out.} \\
\\ 
\texttt{Your job is to infer what it does and then submit a correct Python implementation.} \\
\\ 
\texttt{Function signature:} \\
\texttt{(text: str) {-}> str} \\
\\ 
\texttt{You have 15 turns total.} \\
\\ 
\texttt{I control which valid input{-}output examples you see. You can request them one at a time.} \\
\\ 
\texttt{You must make your best SOLVE attempt no later than the second{-}to{-}last turn.} \\
\texttt{Do not wait until the very end.} \\
\\ 
\texttt{You can use these actions:} \\
\texttt{{-} JSON format described at the end of the prompt} \\
\texttt{{-} SOLVE: ```python ... ```} \\
\\ 
\texttt{Do not add import statements in your guess.} \\
\texttt{Anything you might need is already available, so just write the function body in hypothesis and return the answer.} \\
\\ 
\texttt{Preloaded examples:} \\
\texttt{1. ('a{-}b', 'a{-}b')} \\
\\ 
\texttt{Scoring:} \\
\texttt{{-} correctness is required for any non{-}zero score} \\
\texttt{{-} in interactive modes, solving with fewer information requests is better} \\
\\ 
\texttt{Input formatting reminders:} \\
\texttt{{-} strings must be quoted} \\
\texttt{{-} booleans must be True/False} \\
\texttt{{-} use commas between multiple values} \\
\texttt{{-} do not write any extra action tags in the same message} \\
\\ 
\texttt{Format example:} \\
\texttt{You write:} \\
\texttt{\{'current\_mode': <probing/confirming/don't know; choose whichever best describes your current state in the task>, 'input\_rationale': <brief description of why you chose this hypothesis>, 'hypothesis': <either the exact string "unknown" if you have no guess yet OR your current best guess as runnable Python matching the given signature, with a one{-}line docstring stating what it does. Never write placeholder code, if you cannot commit to a guess, use "unknown". Match the given signature.} \\
\texttt{e.g.} \\
\texttt{def f(x):} \\
\texttt{	"""Returns the absolute value of x."""} \\
\texttt{	return x if x >= 0 else {-}x} \\
\\ 
\texttt{Use "unknown" if you have no guess.>, 'confidence\_score': <score between 1{-}100. 100 means highest confidence in your hypothesis>\}} \\
\\ 
\texttt{Output only one valid json per turn.} \\
            }
        }
    }
     \\ \\

    \theutterance \stepcounter{utterance}  
    & \multicolumn{4}{p{0.6\linewidth}}{
        \cellcolor[rgb]{0.561,0.737,0.561}{
            \makecell[{{p{\linewidth}}}]{
                \texttt{\tiny{[P1$\rangle$GM]}}
                \texttt{```json} \\
\texttt{\{} \\
\texttt{  "current\_mode": "probing",} \\
\texttt{  "input\_rationale": "The preloaded example ('a{-}b', 'a{-}b') suggests that the function might return the input unchanged in simple cases. I will request an example with a more complex input to see if there are any transformations or patterns.",} \\
\texttt{  "hypothesis": "unknown",} \\
\texttt{  "confidence\_score": 1} \\
\texttt{\}} \\
\texttt{```} \\
            }
        }
    }
    & & \\ \\

    \theutterance \stepcounter{utterance}  
    & & & \multicolumn{4}{p{0.6\linewidth}}{
        \cellcolor[rgb]{0.9,0.9,0.9}{
            \makecell[{{p{\linewidth}}}]{
                \texttt{\tiny{[P1$\langle$GM]}}
                \texttt{OUTPUT: ('hello', 'llo')} \\
            }
        }
    }
     \\ \\

    \theutterance \stepcounter{utterance}  
    & \multicolumn{4}{p{0.6\linewidth}}{
        \cellcolor[rgb]{0.561,0.737,0.561}{
            \makecell[{{p{\linewidth}}}]{
                \texttt{\tiny{[P1$\rangle$GM]}}
                \texttt{```json} \\
\texttt{\{} \\
\texttt{  "current\_mode": "probing",} \\
\texttt{  "input\_rationale": "The new example ('hello', 'llo') suggests that the function might be removing the first character of the string. However, the preloaded example ('a{-}b', 'a{-}b') contradicts this, as the output is unchanged. This indicates the function might have conditional behavior, possibly removing a prefix under certain conditions (e.g., if the string starts with a specific pattern). I will request another example to test this hypothesis.",} \\
\texttt{  "hypothesis": "unknown",} \\
\texttt{  "confidence\_score": 30} \\
\texttt{\}} \\
\texttt{```} \\
            }
        }
    }
    & & \\ \\

    \theutterance \stepcounter{utterance}  
    & & & \multicolumn{4}{p{0.6\linewidth}}{
        \cellcolor[rgb]{0.9,0.9,0.9}{
            \makecell[{{p{\linewidth}}}]{
                \texttt{\tiny{[P1$\langle$GM]}}
                \texttt{OUTPUT: ('   ', '   ')} \\
            }
        }
    }
     \\ \\

    \theutterance \stepcounter{utterance}  
    & \multicolumn{4}{p{0.6\linewidth}}{
        \cellcolor[rgb]{0.561,0.737,0.561}{
            \makecell[{{p{\linewidth}}}]{
                \texttt{\tiny{[P1$\rangle$GM]}}
                \texttt{```json} \\
\texttt{\{} \\
\texttt{  "current\_mode": "probing",} \\
\texttt{  "input\_rationale": "The example ('   ', '   ') shows that whitespace{-}only strings are returned unchanged. This suggests the function does not simply strip whitespace or remove the first character unconditionally. The earlier examples ('a{-}b' → 'a{-}b' and 'hello' → 'llo') hint at a pattern where the function might remove a prefix only if it matches a specific condition (e.g., a single non{-}whitespace character followed by a hyphen, or a lowercase letter at the start). I will request another example to test this.",} \\
\texttt{  "hypothesis": "unknown",} \\
\texttt{  "confidence\_score": 50} \\
\texttt{\}} \\
\texttt{``` \textbackslash{}\textbackslash{}} \\
\texttt{Request: Provide an example with input `'x{-}y{-}z'` or `'test{-}123'`. This will help confirm whether the function removes a prefix before a hyphen or follows another rule.} \\
            }
        }
    }
    & & \\ \\

    \theutterance \stepcounter{utterance}  
    & & & \multicolumn{4}{p{0.6\linewidth}}{
        \cellcolor[rgb]{0.9,0.9,0.9}{
            \makecell[{{p{\linewidth}}}]{
                \texttt{\tiny{[P1$\langle$GM]}}
                \texttt{Game Violation: ParseError: Response is not valid JSON.} \\
\texttt{Output ONLY:} \\
\texttt{One JSON with confidence\_score, current\_mode, input\_rationale, and hypothesis keys ...} \\
\texttt{OR} \\
\texttt{SOLVE: ```python ... ```} \\
            }
        }
    }
     \\ \\

    \theutterance \stepcounter{utterance}  
    & \multicolumn{4}{p{0.6\linewidth}}{
        \cellcolor[rgb]{0.561,0.737,0.561}{
            \makecell[{{p{\linewidth}}}]{
                \texttt{\tiny{[P1$\rangle$GM]}}
                \texttt{```json} \\
\texttt{\{} \\
\texttt{  "current\_mode": "probing",} \\
\texttt{  "input\_rationale": "The previous examples suggest the function may remove a prefix up to the first hyphen if certain conditions are met (e.g., 'hello' → 'llo' removes 'he', while 'a{-}b' remains unchanged). The whitespace example shows no transformation for non{-}alphabetic inputs. I will test whether the function removes the first 'word' (alphabetic characters) before a hyphen or another delimiter.",} \\
\texttt{  "hypothesis": "unknown",} \\
\texttt{  "confidence\_score": 60} \\
\texttt{\}} \\
\texttt{``` \textbackslash{}\textbackslash{}} \\
\texttt{Request: Provide the output for input `'test{-}me'`. This will help determine if the function splits on hyphens or follows another rule.} \\
            }
        }
    }
    & & \\ \\

    \theutterance \stepcounter{utterance}  
    & & & \multicolumn{4}{p{0.6\linewidth}}{
        \cellcolor[rgb]{0.9,0.9,0.9}{
            \makecell[{{p{\linewidth}}}]{
                \texttt{\tiny{[P1$\langle$GM]}}
                \texttt{Game Violation: ParseError: Response is not valid JSON.} \\
\texttt{Output ONLY:} \\
\texttt{One JSON with confidence\_score, current\_mode, input\_rationale, and hypothesis keys ...} \\
\texttt{OR} \\
\texttt{SOLVE: ```python ... ```} \\
            }
        }
    }
     \\ \\

    \theutterance \stepcounter{utterance}  
    & \multicolumn{4}{p{0.6\linewidth}}{
        \cellcolor[rgb]{0.561,0.737,0.561}{
            \makecell[{{p{\linewidth}}}]{
                \texttt{\tiny{[P1$\rangle$GM]}}
                \texttt{```json} \\
\texttt{\{} \\
\texttt{  "current\_mode": "probing",} \\
\texttt{  "input\_rationale": "Based on the examples so far ('a{-}b' → 'a{-}b', 'hello' → 'llo', '   ' → '   '), the function appears to remove the first alphabetic segment (if it exists) before a hyphen or the rest of the string if no hyphen is present. To confirm, I will request an example with a hyphenated string like 'test{-}me' or 'x{-}y{-}z'.",} \\
\texttt{  "hypothesis": "def f(text: str) {-}> str:\textbackslash{}n    \textbackslash{}"\textbackslash{}"\textbackslash{}"Removes the first alphabetic segment before a hyphen or the entire prefix if no hyphen exists.\textbackslash{}"\textbackslash{}"\textbackslash{}"\textbackslash{}n    if '{-}' in text:\textbackslash{}n        parts = text.split('{-}', 1)\textbackslash{}n        if parts{[}0{]}.isalpha():\textbackslash{}n            return parts{[}1{]} if parts{[}1{]} else text\textbackslash{}n    return text.lstrip('abcdefghijklmnopqrstuvwxyz') or text",} \\
\texttt{  "confidence\_score": 70} \\
\texttt{\}} \\
\texttt{``` \textbackslash{}\textbackslash{}} \\
\texttt{Request: Provide the output for input `'test{-}me'`.} \\
            }
        }
    }
    & & \\ \\

    \theutterance \stepcounter{utterance}  
    & & & \multicolumn{4}{p{0.6\linewidth}}{
        \cellcolor[rgb]{0.9,0.9,0.9}{
            \makecell[{{p{\linewidth}}}]{
                \texttt{\tiny{[P1$\langle$GM]}}
                \texttt{Game Violation: ParseError: Response is not valid JSON.} \\
\texttt{Output ONLY:} \\
\texttt{One JSON with confidence\_score, current\_mode, input\_rationale, and hypothesis keys ...} \\
\texttt{OR} \\
\texttt{SOLVE: ```python ... ```} \\
            }
        }
    }
     \\ \\

    \theutterance \stepcounter{utterance}  
    & \multicolumn{4}{p{0.6\linewidth}}{
        \cellcolor[rgb]{0.561,0.737,0.561}{
            \makecell[{{p{\linewidth}}}]{
                \texttt{\tiny{[P1$\rangle$GM]}}
                \texttt{```json} \\
\texttt{\{} \\
\texttt{  "current\_mode": "confirming",} \\
\texttt{  "input\_rationale": "Based on the observed behavior, the function likely removes the first alphabetic segment before a hyphen (if present) or the first lowercase alphabetic prefix (if no hyphen exists). The preloaded example ('a{-}b' → 'a{-}b') suggests it preserves the string if the prefix is a single character, while ('hello' → 'llo') shows removal of the first lowercase segment. I will now submit a SOLVE attempt.",} \\
\texttt{  "hypothesis": "def f(text: str) {-}> str:\textbackslash{}n    \textbackslash{}"\textbackslash{}"\textbackslash{}"Removes the first lowercase alphabetic segment before a hyphen or the first lowercase prefix if no hyphen exists.\textbackslash{}"\textbackslash{}"\textbackslash{}"\textbackslash{}n    if '{-}' in text:\textbackslash{}n        prefix, suffix = text.split('{-}', 1)\textbackslash{}n        if prefix.islower() and prefix.isalpha():\textbackslash{}n            return suffix if suffix else text\textbackslash{}n    return text.lstrip('abcdefghijklmnopqrstuvwxyz') or text",} \\
\texttt{  "confidence\_score": 85} \\
\texttt{\}} \\
\texttt{``` \textbackslash{}\textbackslash{}} \\
\texttt{SOLVE:} \\
\texttt{```python} \\
\texttt{def f(text: str) {-}> str:} \\
\texttt{    """Removes the first lowercase alphabetic segment before a hyphen or the first lowercase prefix if no hyphen exists."""} \\
\texttt{    if '{-}' in text:} \\
\texttt{        prefix, suffix = text.split('{-}', 1)} \\
\texttt{        if prefix.islower() and prefix} \\
            }
        }
    }
    & & \\ \\

    \theutterance \stepcounter{utterance}  
    & & & \multicolumn{4}{p{0.6\linewidth}}{
        \cellcolor[rgb]{0.9,0.9,0.9}{
            \makecell[{{p{\linewidth}}}]{
                \texttt{\tiny{[P1$\langle$GM]}}
                \texttt{Game Violation: ParseError: Response is not valid JSON.} \\
\texttt{Output ONLY:} \\
\texttt{One JSON with confidence\_score, current\_mode, input\_rationale, and hypothesis keys ...} \\
\texttt{OR} \\
\texttt{SOLVE: ```python ... ```} \\
            }
        }
    }
     \\ \\

    \theutterance \stepcounter{utterance}  
    & \multicolumn{4}{p{0.6\linewidth}}{
        \cellcolor[rgb]{0.561,0.737,0.561}{
            \makecell[{{p{\linewidth}}}]{
                \texttt{\tiny{[P1$\rangle$GM]}}
                \texttt{SOLVE:} \\
\texttt{```python} \\
\texttt{def f(text: str) {-}> str:} \\
\texttt{    """Returns the string after removing the first lowercase alphabetic prefix before a hyphen (if present) or the first lowercase alphabetic prefix (if no hyphen exists)."""} \\
\texttt{    if '{-}' in text:} \\
\texttt{        prefix, suffix = text.split('{-}', 1)} \\
\texttt{        if prefix.islower() and prefix.isalpha():} \\
\texttt{            return suffix} \\
\texttt{    stripped = text.lstrip('abcdefghijklmnopqrstuvwxyz')} \\
\texttt{    return stripped if stripped else text} \\
\texttt{```} \\
            }
        }
    }
    & & \\ \\

    \theutterance \stepcounter{utterance}  
    & & & \multicolumn{2}{p{0.3\linewidth}}{
        \cellcolor[rgb]{0.9,0.9,0.9}{
            \makecell[{{p{\linewidth}}}]{
                \texttt{\tiny{[GM$|$GM]}}
                \texttt{Game Verdict: LOSS} \\
            }
        }
    }
    & & \\ \\

    \theutterance \stepcounter{utterance}  
    & & & \multicolumn{2}{p{0.3\linewidth}}{
        \cellcolor[rgb]{0.9,0.9,0.9}{
            \makecell[{{p{\linewidth}}}]{
                \texttt{\tiny{[GM$|$GM]}}
                \texttt{Incorrect.} \\
\texttt{Failed Cases:} \\
\texttt{Input: {[}'a{-}b'{]} | Expected: a{-}b | Got: b} \\
\texttt{Input: {[}'hello'{]} | Expected: llo | Got: hello} \\
\texttt{Input: {[}'hello world'{]} | Expected: rld | Got:  world} \\
            }
        }
    }
    & & \\ \\

    \theutterance \stepcounter{utterance}  
    & & & \multicolumn{2}{p{0.3\linewidth}}{
        \cellcolor[rgb]{0.9,0.9,0.9}{
            \makecell[{{p{\linewidth}}}]{
                \texttt{\tiny{[GM$|$GM]}}
                \texttt{def take\_last\_three(text: str) {-}> str:} \\
\texttt{    return text{[}{-}3:{]}} \\
            }
        }
    }
    & & \\ \\

    \theutterance \stepcounter{utterance}  
    & & & \multicolumn{2}{p{0.3\linewidth}}{
        \cellcolor[rgb]{0.9,0.9,0.9}{
            \makecell[{{p{\linewidth}}}]{
                \texttt{\tiny{[GM$|$GM]}}
                \texttt{def f(text: str) {-}> str:} \\
\texttt{    """Returns the string after removing the first lowercase alphabetic prefix before a hyphen (if present) or the first lowercase alphabetic prefix (if no hyphen exists)."""} \\
\texttt{    if '{-}' in text:} \\
\texttt{        prefix, suffix = text.split('{-}', 1)} \\
\texttt{        if prefix.islower() and prefix.isalpha():} \\
\texttt{            return suffix} \\
\texttt{    stripped = text.lstrip('abcdefghijklmnopqrstuvwxyz')} \\
\texttt{    return stripped if stripped else text} \\
            }
        }
    }
    & & \\ \\

    \theutterance \stepcounter{utterance}  
    & & & \multicolumn{2}{p{0.3\linewidth}}{
        \cellcolor[rgb]{0.9,0.9,0.9}{
            \makecell[{{p{\linewidth}}}]{
                \texttt{\tiny{[GM$|$GM]}}
                \texttt{Mode: passive\_examples} \\
\texttt{Accuracy: 0.220} \\
\texttt{Information requests: 2/15} \\
\texttt{Revealed passive examples: 3} \\
\texttt{Efficiency: 0.933} \\
\texttt{Slack: 13} \\
\texttt{Observed{-}pair consistency: 0.333} \\
\texttt{Parse errors: 4} \\
\texttt{Runtime errors: 0} \\
            }
        }
    }
    & & \\ \\

\end{supertabular}
}

\end{document}
